\documentclass{beamer}
%\documentclass[handout]{beamer}

\mode<presentation>
{
%  \usetheme{Warsaw}
  % or ...

%  \setbeamercovered{transparent}
  % or whatever (possibly just delete it)
}


\setbeamertemplate{navigation symbols}{}

\usepackage[utf8]{inputenc}
\usepackage[spanish, es-tabla, es-nodecimaldot]{babel}
\usepackage{xcolor}

\usepackage{listings}

%New colors defined below
\definecolor{codegreen}{rgb}{0,0.6,0}
\definecolor{codegray}{rgb}{0.5,0.5,0.5}
\definecolor{codepurple}{rgb}{0.58,0,0.82}
\definecolor{backcolour}{rgb}{0.95,0.95,0.92}

%Code listing style named "mystyle"
\lstdefinestyle{mystyle}{
  backgroundcolor=\color{backcolour},   commentstyle=\color{codegreen},
  keywordstyle=\color{magenta},
  numberstyle=\tiny\color{codegray},
  stringstyle=\color{codepurple},
  basicstyle=\ttfamily\footnotesize,
  breakatwhitespace=false,         
  breaklines=true,                 
  captionpos=b,                    
  keepspaces=true,                 
  numbers=left,                    
  numbersep=5pt,                  
  showspaces=false,                
  showstringspaces=false,
  showtabs=false,                  
  tabsize=2
}

%"mystyle" code listing set
\lstset{style=mystyle}


%\usepackage{times}
%\usepackage[T1]{fontenc}
% Or whatever. Note that the encoding and the font should match. If T1
% does not look nice, try deleting the line with the fontenc.


\title[Intro Prog] % (optional, use only with long paper titles)
{Introducción a Python}

\subtitle
{Clase 3}

\author[MS]
{Mateo Suster \\ mateosuster@gmail.com}%\inst{1}}
% - Give the names in the same order as the appear in the paper.
% - Use the \inst{?} command only if the authors have different
%   affiliation.

\institute[UNGS] % (optional, but mostly needed)
{
%  \inst{1}%
  Matemática para Economistas III \\ 
  Instituto de Industria\\
  Universidad Nacional de General Sarmiento
}

\date[] % (optional, should be abbreviation of conference name)
{ \today}



\begin{document}

\begin{frame}
  \titlepage
\end{frame}


% PAG 1
\begin{frame}[fragile]{¿Qué hace el siguiente programa en Python?} 

\begin{lstlisting}[language=Python]
#ingreso las millas 
mills = 34.122

#hago la conversion a kilometros
km = mills * 1.6

km_redondo = round(km, 2)

#imprimo por pantalla el resultado
print(mills, "millas son", km_redondo, "kilometros")

#la salida queda 
> 34.122 millas son 54.59 kilometros
\end{lstlisting} \pause

\begin{itemize}
    \item Un programa sencillo de conversión, que a esta altura nos debe resultar familiar
\end{itemize}
\end{frame}


%PAG 2
\begin{frame}[fragile]{El mismo programa, pero con listas}
\begin{lstlisting}[language=Python]
#ingreso las millas 
mills = [34.122, 17.588, 3.187] 

# Hago la conversion a kilometros
km = [mills[0]*1.6, mills[1]*1.6, mills[2]*1.6]

# Imprimo por pantalla el resultado
print(mills[0],"millas son",km[0],"kilometros")  
print(mills[1],"millas son",km[1],"kilometros")  
print(mills[2],"millas son",km[2],"kilometros")
\end{lstlisting} \pause

La salida:
\begin{lstlisting}
> 34.122 millas son 54.5952 kilometros
> 17.588 millas son 28.1408 kilometros
> 3.187 millas son 5.0992 kilometros
\end{lstlisting} \pause

\begin{itemize}
    \item ¿Qué problemas tiene esta codificación? \pause Información duplicada y \textbf{sobretodo ineficiencia} para programar un proceso automatizable por medio de una \textbf{estructura de control}. \pause
    \item Tener más codigo no es síntoma de ser un buen programador
\end{itemize}
\end{frame}


%PAG 3
\begin{frame}[fragile]{El mismo programa, pero con ciclo while}
\begin{lstlisting}[language=Python]
#ingreso las millas 
mills = [1.1, 33.4, 34.122, 17.588, 3.187, 50.] 
    
#guardo la longitud de la lista
longitud = len(mills)
    
#hago la conversion a kilometros
km = [0]*longitud #lista de ceros de longitud adecuada
    
i = 0 # variable para indicar la posicion en la lista
while i < longitud: #condicion
    km[i] = mills[i]*1.6 #conversion
    i = i +1 #avance de posicion
print(km)
#la salida queda igual a: 
> [1.8, 53.4, 54.6, 28.1, 5.1, 80.0] 
\end{lstlisting} \pause

\begin{itemize}
        \item Un programa automatizado de conversión, con el que se tienen que familiarizar. \pause
        \item ¿Pero se puede automatizar aún más? \pause Sí.
\end{itemize}
\end{frame}


% PAG 4
\begin{frame}{Funciones} \pause
\begin{itemize}
    \item Una función es una unidad de código que aísla una parte de un cómputo. \textbf{Es un programa dentro de un programa.} \pause
    \item Poseen argumentos o parámetros (inputs), los procesa y devuelve una salida (output). \pause (Piensen en las funciones matemáticas) \pause
    \item Características de las funciones: \pause
\end{itemize}
\begin{enumerate}
    \item Permite \textbf{dividir} un problema en problemas más simples. \pause
    \item Permite \textbf{ordenar conceptualmente} el código para que sea más fácil de entender. \pause
    \item Permite \textbf{reutilizar soluciones} a problemas pequeños en la solución de problemas mayores. 
\end{enumerate}
\end{frame}


%PAG 5
\begin{frame}[fragile]{Estructura de las funciones} \pause


\begin{itemize}
	\item Funciones que retornan un valor
\end{itemize}
\begin{lstlisting}[language=Python]
def NOMBRE(ARGUMENTOS):
    ENUNCIADOS 
    return VALOR
\end{lstlisting} \pause

\begin{itemize}
        \item Funciones que no devuelven valor alguno (también llamadas Procedimiento)   
\end{itemize}
\begin{lstlisting}[language=Python]
def NOMBRE(ARGUMENTOS):
    ENUNCIADOS
    print(VALOR)
\end{lstlisting} \pause

\begin{itemize}
	\item Nuevamente, atención con el bloque indentado! \pause
	\item Tanto en los enunciados (donde se realizan las operaciones sobre los argumentos), como en la salida, puede ir cualquier cosa. \pause
	\item \alert{Importante:} \texttt{def} y \texttt{return} son dos palabras reservadas. \pause Se puede retornar más de un valor
\end{itemize}

\end{frame}


% MEME FUNCION
\begin{frame}{En los enunciados puede ir cualquier cosa cualquier cosa...} \pause
\begin{center}
\includegraphics[height=7cm, scale=0.5]{meme_funciones.png}
\end{center}
\end{frame}






%PAG 6
\begin{frame}[fragile]{El mismo programa, pero en una función} \pause

\begin{lstlisting}[language=Python]
def convertir(lista):
    longitud = len(lista) 
   # Hago la conversion
    listaConvertida = [0]*longitud # lista de ceros
    i = 0
    while i < longitud:
        listaConvertida[i] = lista[i]*1.6
        i = i + 1
    return listaConvertida

# El resto del codigo queda mas prolijo
mills = [1.1, 33.4, 13.2, 60.0, 17.35]
kms = convertir(mills)
print(kms)
> [1.8, 53.4, 54.6, 28.1, 5.1, 80.0] #  la salida
\end{lstlisting} \pause

\begin{itemize}
    \item La función \texttt{convertir} \textbf{encapsula} el problema de la conversión a kilómetros. \pause
    \item Ventajas? \pause \textbf{Particionamos} el problema en otros dos de menor complejidad y el programa \textbf{puede ser resuelto por diferentes personas y al mismo tiempo}
\end{itemize}
\end{frame}



%PAG 7
\begin{frame}[fragile]{¿Cuál es la diferencia con el programa anterior?} \pause

\begin{lstlisting}[language=Python]
def convertir(lista,factor):
    longitud = len(lista) #guardo la longitud de la lista
   # Hago la conversion
    listaConvertida = [0]*longitud # lista de ceros
    i = 0
    while i < longitud:
        listaConvertida[i] = lista[i]*factor
        i = i + 1
    return listaConvertida

# El resto del codigo queda mas prolijo
mills = [1.1, 33.4, 13.2, 60.0, 17.35]
fc = 1.6
kms = convertir(mills,fc)
\end{lstlisting} \pause

\begin{itemize}
    \item La función convertir ahora toma dos parámetros de entrada. \pause
    \item La ventaja de esto es que ahora la función \texttt{convertir} funciona para cualquier lista y para cualquier factor de conversión.
\end{itemize}
\end{frame}






% PAG final
\begin{frame}{A coedear se a dicho!} \pause
Nuestro nuevo aprendizaje:  \pause
\begin{itemize}
	\item Siempre que sea posible \alert{particionar} el problema y \alert{encapsularlo} en problemas menos complejos mediante \alert{funciones}
\end{itemize} 
\end{frame}





\end{document}
